
%\thispagestyle{empty}

%\printnomenclature
\renewcommand\nomgroup[1]{%
  \item[\bfseries
  \ifstrequal{#1}{F}{Fluid properties}{%
  \ifstrequal{#1}{O}{Other symbols}{}}%
]}
%%%%%%%%%%%%%%%%%%%%%%%%%%%%%%%%%%%%%%%%%%%%%%%%%%%%%%%%%%%%%%
% Wenn Symbolverzeichnis automatisch erzeugt werdenn soll

\ifnomenklatur
%\ifenglish
%\printnomenclature\addcontentsline{toc}{chapter}{Notation}
%\else
%\printnomenclature\addcontentsline{toc}{chapter}{Nomenklatur}
%\fi

%%%%%%%%%%%%%%%%%%%%%%%%%%%%%%%%%%%%%%%%%%%%%%%%%%%%%%%%%%%%%%%%
\else
%% Wenn man das Symbolverzeichnis selber machen will
%\ifenglish
%\chapter*{Notation}\addcontentsline{toc}{chapter}{Notation}
%\else
\chapter*{List of Symbols}\addcontentsline{toc}{chapter}{List of Symbols}
%\fi
%
%
%
%
%%\thispagestyle{empty}
%\ifenglish
%\section*{Abbrevations}
%\else
%\section*{Abkürzungen}
%\fi
%
\nomenclature[F]{$\mu$}{Dynamic viscosity of the fluid}
\nomenclature[F]{$\rho$}{Density of the fluid}
\nomenclature[O]{$Re_b$}{Bulk Reynolds Number}
\nomenclature[O]{$Re_\tau$}{Friction Reynolds number}
\nomenclature[O]{C}{Mass flux}
\nomenclature[O]{$\alpha$}{Scaling factor to ensure constant mass flux}


\printnomenclature

% \begin{longtable}[t]{p{0.1\textwidth}p{0.1\textwidth}p{0.8\textwidth}}
% 	\textbf{Variable}  & \textbf{Einheit} & \textbf{Beschreibung}\\
% 	\midrule
% 	$c_0$ & $m/s$ & Schallgeschwindigkeit in Luft \\
% 	$\nabla$ & $1/m$ & Nabla-Operator \\
% 	$\mu$ & & Dynamic viscosity of the fluid\\
% 	$\rho$ & & Density of the fluid\\
% 	$\alpha$ & & Scaling factor to ensure constant mass flux\\
% 	$Re_\tau$ & & Friction Reynolds number\\
% 	$Re_b$ & & Bulk Reynolds number\\
% 	$C$ & & Mass flux\\

	
% 	%
% \end{longtable}


%\chapter*{Abkürzungsverzeichnis}\addcontentsline{toc}{chapter}{Abkürzungsverzeichnis}
%{\let\clearpage\relax \chapter*{List of Abbreviations}}\addcontentsline{toc}{chapter}{List of Abbreviations}

\newacronym{sem}{SEM}{Spectral Element}
\newacronym{cfd}{CFD}{computational fluid dynamics}
\newacronym{dns}{DNS}{Direct Numerical Simulation}
\newacronym{les}{LES}{Large Eddy Simulation}
\newacronym{nse}{NSE}{Navier-Stokes-Equations}
\newacronym{rans}{RANS}{Reynolds-averaged Navier-Stokes}
\newacronym{da}{DA}{Double-Averaged}
\newacronym{bc}{B.C}{Boundary Condition}
\newacronym{bcs}{B.Cs}{Boundary Conditions}
\newacronym{ic}{I.C}{Initial Condition}
\newacronym{ics}{I.Cs}{Initial Conditions}
\newacronym{gll}{GLL}{Gauss-Lobatto-Legrende}
\newacronym{fem}{FEM}{Finite-Element-Method}
\printglossary[type=\acronymtype]


% \begin{longtable}[t]{p{0.25\textwidth}p{0.8\textwidth}}
% 	\textbf{Abkürzung}  & \textbf{Bedeutung}\\
% 	\midrule
% 	NSE & Navier-Stokes Equations \\
	

	
	
% 	%
% \end{longtable}



\fi	
\cleardoublepage